
[[1.1]]
Exercise 1: Natural units (5 points)
Natural units are defined by setting $c=\hbar=1$.
Use natural units to express 1 kg in GeV , as well as 1 s in $1 / \mathrm{GeV}$. Use your results to express Newton's constant of gravity

$$
G_N=6,67 \times 10^{-11} \frac{\mathrm{~m}^3}{\mathrm{~kg} \mathrm{~s}^2}
$$

in natural units.
Finally give the value of the Planck mass $M_{p l}=1 / \sqrt{G_N}$.

[[2.1]]
Exercise 2: Charged particle in a constant magnetic field (5 points)
We consider a particle with charge $e$ and mass $m$. The Lagrangian of a charged particle in an electromagnetic field is given by

$$
L(\mathbf{r}, \dot{\mathbf{r}}, t)=\frac{1}{2} m \dot{\mathbf{r}}^2+e \frac{\dot{\mathbf{r}}}{c} \cdot \mathbf{A}(\mathbf{r})-e \Phi(\mathbf{r})
$$

with the scalar potential $\Phi(\mathbf{r})$ and the vector potential $\mathbf{A}(\mathbf{r})$.
Calculate the Hamilton function $H(\mathbf{r}, \mathbf{p}, t)$ of the system, using a Legendre transformation.

[[3.1]]
Exercise 3: Continuity equation in quantum mechanics (3 points)
Consider a wave function $\psi(\mathbf{r}, t)$ satisfying the Schrödinger equation

$$
\left(-\frac{\hbar^2}{2 m} \Delta+V(\mathbf{r})\right) \psi(\mathbf{r}, t)=\mathrm{i} \hbar \frac{\partial}{\partial t} \psi(\mathbf{r}, t) .
$$


The probability density is defined as $\rho(\mathbf{r}, t):=\psi^*(\mathbf{r}, t) \psi(\mathbf{r}, t)$. Show that it satisfies the continuity equation

$$
\frac{\partial}{\partial t} \rho(\mathbf{r}, t)+\nabla \cdot \mathbf{j}(\mathbf{r}, t)=0,
$$

where $\mathbf{j}(\mathbf{r}, t)$ is the propability current

$$
\mathbf{j}(\mathbf{r}, t):=\frac{\hbar}{2 \mathrm{i} m}\left(\psi^*(\nabla \psi)-\left(\nabla \psi^*\right) \psi\right) .
$$

[[4.1]]
Exercise 4: Harmonic oscillator in quantum mechanics (7 points)
Recall the time-independent Schrödinger equation of the one-dimensional harmonic oscillator in quantum mechanics

$$
\left(\frac{\hat{p}^2}{2 m}+\frac{1}{2} m \omega^2 \hat{x}^2\right) \psi(x)=E \psi(x),
$$

with momentum operator $\hat{p}$ and the position operator $\hat{x}$.
It has been shown that the energy eigenvalues $E$ can be calculated easily in an energy basis of the system $\{|n\rangle\}, n \in \mathbb{N}$, by making use of two operators defined as

$$
\hat{a}:=\sqrt{\frac{m \omega}{2 \hbar}}\left(\hat{x}+\frac{\mathrm{i}}{m \omega} \hat{p}\right), \quad \hat{a}^{\dagger}:=\sqrt{\frac{m \omega}{2 \hbar}}\left(\hat{x}-\frac{\mathrm{i}}{m \omega} \hat{p}\right),
$$

with the property to raise and lower ('create' and 'annihilate') the states

$$
\hat{a}|n\rangle=c_n^{-}|n-1\rangle, \quad \hat{a}^{\dagger}|n\rangle=c_n^{+}|n+1\rangle
$$


The occupation number operator acts as

$$
\hat{N}|n\rangle=\hat{a}^{\dagger} \hat{a}|n\rangle=n|n\rangle .
$$


Calculate the commutator $\left[\hat{a}, \hat{a}^{\dagger}\right]$ and the normalization constants $c_n^{-}$and $c_n^{+}$. Use the operators $\hat{a}$ and $\hat{a}^{\dagger}$ to reformulate the Hamilton operator and give the energy eigenvalues $E_n$.